% Generated draft: requires mathematical and source review.
\documentclass[11pt]{article}
\usepackage[margin=25mm]{geometry}
\usepackage{fontspec}
\setmainfont{Latin Modern Roman}
\usepackage{amsmath,amssymb,mathtools,unicode-math}
\setmathfont{Latin Modern Math}
\usepackage{enumitem,needspace,xurl,hyperref}
\setlength{\parindent}{0pt}
\setlength{\parskip}{5pt}
\setlength{\emergencystretch}{3em}
\newcommand{\N}{\mathbb{N}}
\newcommand{\Z}{\mathbb{Z}}
\newcommand{\Q}{\mathbb{Q}}
\newcommand{\R}{\mathbb{R}}
\newcommand{\C}{\mathbb{C}}
\newcommand{\F}{\mathbb{F}}
\newcommand{\A}{\mathbb{A}}
\newcommand{\PP}{\mathbb{P}}
\newcommand{\E}{\mathbb{E}}
\newcommand{\eps}{\varepsilon}
\newcommand{\dd}{\,\mathrm{d}}
\begin{document}
\section*{UnsolvedMath ID: 20002816}
\textbf{Source problem code:} catalogue entry
\subsection{Definitions and mathematical statement}
Problem 1.4 asks to characterize the Andr\'e--Yan supercharacter theory ``along the lines of the previous problem''. The preceding local registry record, permanent ID 20002815, asks whether Andr\'e's theory is the coarsest supercharacter theory containing the single-box elementary irreducibles individually. This supplies a natural minimality interpretation, rather than an unspecified characterization. The neighboring record is context and is not altered or completed here.

Fix a prime power \(q\), \(n\ge1\), \(G=U_n(q)\), and the algebra \(J\) of strictly upper-triangular matrices over \(\F_q\), so \(G=1+J\). Fix a nontrivial additive character \(\psi:\F_q\to\C^*\). In the standard algebra-group formulation the Andr\'e--Yan superclasses are
\[
 1+G XG\quad(X\in J).
\]
For \(\lambda\in J^*\), the associated supercharacter can be normalized as
\[
 \chi_\lambda(g)=\frac{|G\lambda|}{|G\lambda G|}
               \sum_{\nu\in G\lambda G}\psi(\nu(g-1)).
\]
The dual left/right actions are those induced by multiplication on \(J\). The single-box elementary characters are \(\xi_{ij}(a)=\chi_{\lambda_{ij,a}}\), where \(i<j\), \(a\in\F_q^*\), and \(\lambda_{ij,a}(X)=aX_{ij}\). The aid uses the known irreducibility of these elementary characters. For \(n=1\) the set is empty and the group is trivial.

A supercharacter theory consists of a partition of \(G\) into superclasses and a compatible partition of \(\operatorname{Irr}(G)\), with singleton identity superclass and the usual equal-size and character-constancy conditions. Retaining an elementary irreducible individually means that its singleton is a block of the irreducible-character partition, not merely that it occurs somewhere in a reducible supercharacter. Coarser theories have coarser superclass partitions.

\textbf{Characterization reported in the supplied aid.} Let \(\operatorname{SC}_{\mathrm{AY}}(G)\) be the complex algebra of all functions constant on Andr\'e--Yan superclasses, with pointwise multiplication. For every \(n\) and prime power \(q\), the aid reports
\[
 \operatorname{SC}_{\mathrm{AY}}(G)
       =\C[\xi_{ij}(a):i<j,\ a\in\F_q^*],
\]
where the right side is the unital pointwise function algebra. Equivalently, for \(g,h\in G\),
\[
 g,h\text{ lie in the same Andr\'e--Yan superclass}
 \quad\Longleftrightarrow\quad
 \xi_{ij}(a)(g)=\xi_{ij}(a)(h)
       \text{ for every }i<j,\ a\in\F_q^*.
\]
Consequently the Andr\'e--Yan theory is the unique coarsest supercharacter theory retaining all these elementary irreducibles as singleton character blocks. This is a characterization relative to a distinguished prescribed family, not a classification of all supercharacter theories of \(U_n(q)\).

The aid attributes the underlying input to Andr\'e's arbitrary-characteristic factorization and unique-support results: basic supercharacters factor as products of elementary single-box characters, and irreducible supports form the required partition. The finite-function-algebra interpretation then identifies superclasses with simultaneous value fibers. Full authoritative references and normalization details for those inputs are not supplied. The assertion is a reported supplied result, while the registry remains partially solved and the attempt's problem-status field says open; neither metadata label is independently resolved here.
\subsection{Short English statement}
Using the preceding problem's minimality question, the aid characterizes the Andr\'e–Yan theory by its single-box elementary characters. They generate its superclass-function algebra, so their simultaneous value fibers are exactly the superclasses. Equivalently it is the coarsest theory keeping those irreducibles individually, rather than a classification of every supercharacter theory.
\subsection{Sources}
\begin{itemize}
\item Local UnsolvedMath registry, permanent ID 20002816, catalogue entry; canonical \texttt{aim-representation-theory-notes.json}, notes[87]. Neighboring local record 20002815, AIM-REPRESENTATION\_THEORY-0087, supplies Problem 1.3's single-box/coarsest-theory context.
\item AIM \emph{Supercharacters and combinatorial Hopf algebras}, Problem 1.4, \url{https://aimath.org/WWN/supercharacters/supercharacters.pdf}. The aid reports agreement with the original PDF; this draft does not independently check it. No individual proposer is supplied.
\item Andr\'e and Yan are named by the source theory; the aid invokes Andr\'e's arbitrary-characteristic factorization and unique-support theorem without full bibliographic entries.
\item Machine-generated research attempt 1, partial result, attempt status open, reported rigor proved and low novelty confidence; registry status partially solved. Its generator/fiber/minimality assertion is not peer reviewed and requires independent verification before citation.
\end{itemize}
\end{document}
